% Preamble for main.tex

% ---- Fonts and encoding ----------------------------------------------------
\usepackage[T1]{fontenc}
\usepackage[utf8]{inputenc}
\usepackage{lmodern}
\usepackage{microtype}

% ---- Page layout -----------------------------------------------------------
\usepackage[a4paper, margin=2.5cm, headheight=14pt]{geometry}
\usepackage{fancyhdr}
% Space after the chapter title = space between "Chapter N" and the title (20pt).
\usepackage{etoolbox}
\makeatletter
\patchcmd{\@makechapterhead}{\vskip 40\p@}{\vskip 20\p@}{}{}
\patchcmd{\@makeschapterhead}{\vskip 40\p@}{\vskip 20\p@}{}{}
\makeatother
\pagestyle{fancy}
\fancyhf{}
\fancyhead[LE]{\small\nouppercase{\leftmark}}
\fancyhead[RO]{\small\nouppercase{\rightmark}}
\fancyfoot[C]{\thepage}
\renewcommand{\headrulewidth}{0.4pt}

\setlength{\parindent}{0pt}
\setlength{\parskip}{0.6em}
\raggedbottom

% ---- Mathematics -----------------------------------------------------------
\usepackage{amsmath, amssymb, amsthm}
\usepackage{mathtools}
\usepackage{bm}

% ---- Lists, tables, figures ------------------------------------------------
\usepackage{enumitem}
\setlist{itemsep=0.2em, topsep=0.3em}
\usepackage{booktabs}
\usepackage{graphicx}
\usepackage{xcolor}
\usepackage{pgfplots}
\pgfplotsset{compat=1.18}
\usetikzlibrary{positioning, arrows.meta, patterns, calc, shapes.geometric, decorations.pathreplacing}
\usepgfplotslibrary{fillbetween, groupplots}
% Shared figure styles (see STYLE.md: every image in the notes is redrawn in TikZ/pgfplots)
\pgfplotsset{
  bookplot/.style={width=0.72\textwidth, height=5cm, axis lines=left,
    tick label style={font=\small}, label style={font=\small}, legend style={font=\small, draw=none},
    every axis plot/.append style={thick}, enlarge x limits=0.03,
    scaled y ticks=false, yticklabel style={/pgf/number format/fixed}},
  pmfplot/.style={bookplot, ycomb, ymin=0, every axis plot/.append style={mark=*, mark size=1.5pt, color=thmcolor}},
}
\tikzset{
  venn/.style={draw, thick, fill opacity=0.15},
  block/.style={draw, thick, minimum width=7mm, minimum height=6mm, font=\small},
  dot/.style={circle, fill, inner sep=1.2pt},
}
\usepackage{caption}
\captionsetup{font=small, labelfont=bf}

% ---- Colours (single place to restyle the book) ----------------------------
\definecolor{thmcolor}{HTML}{1F4E79}   % theorems, lemmas, propositions, corollaries
\definecolor{defcolor}{HTML}{2E7D32}   % definitions
\definecolor{excolor}{HTML}{8A6D3B}    % examples
\definecolor{remcolor}{HTML}{6B6B6B}   % remarks
\definecolor{codebg}{HTML}{F6F8FA}
\definecolor{codeframe}{HTML}{D0D7DE}

% ---- Hyperlinks and cross-references ---------------------------------------
\usepackage[colorlinks=true, linkcolor=thmcolor, citecolor=thmcolor, urlcolor=thmcolor]{hyperref}
\usepackage[capitalise, nameinlink]{cleveref}
% (Loaded before the theorem definitions so shared counters get the right names.)
\crefname{example}{Example}{Examples}
\crefname{remark}{Remark}{Remarks}
\crefname{definition}{Definition}{Definitions}

% ---- Theorem-like environments ---------------------------------------------
% One shared counter per chapter (Definition 4.1, Theorem 4.2, Example 4.3, ...)
% so any numbered item is easy to find by number alone.
\theoremstyle{plain}
\newtheorem{theorem}{Theorem}[chapter]
\newtheorem{lemma}[theorem]{Lemma}
\newtheorem{proposition}[theorem]{Proposition}
\newtheorem{corollary}[theorem]{Corollary}

\theoremstyle{definition}
\newtheorem{definition}[theorem]{Definition}
\newtheorem{example}[theorem]{Example}

\theoremstyle{remark}
\newtheorem{remark}[theorem]{Remark}

% Boxes around the environments (amsthm counters keep cleveref working).
\usepackage[most]{tcolorbox}
\tcbset{bookbox/.style={breakable, enhanced, boxrule=0pt, arc=0pt,
  left=6pt, right=6pt, top=4pt, bottom=4pt, before skip=10pt, after skip=10pt}}
\tcolorboxenvironment{theorem}{bookbox, colback=thmcolor!5, borderline west={2.5pt}{0pt}{thmcolor}}
\tcolorboxenvironment{lemma}{bookbox, colback=thmcolor!5, borderline west={2.5pt}{0pt}{thmcolor}}
\tcolorboxenvironment{proposition}{bookbox, colback=thmcolor!5, borderline west={2.5pt}{0pt}{thmcolor}}
\tcolorboxenvironment{corollary}{bookbox, colback=thmcolor!5, borderline west={2.5pt}{0pt}{thmcolor}}
\tcolorboxenvironment{definition}{bookbox, colback=defcolor!5, borderline west={2.5pt}{0pt}{defcolor}}
\tcolorboxenvironment{example}{bookbox, colback=white, borderline west={1pt}{0pt}{excolor}}
\tcolorboxenvironment{remark}{bookbox, colback=white, borderline west={1pt}{0pt}{remcolor}}

% Solution to an example (unnumbered, ends with a marker).
\newenvironment{solution}{\par\textit{Solution.}\ }{\hfill$\blacktriangleleft$\par}

% ---- Code listings ---------------------------------------------------------
\usepackage{listings}
\lstset{
  basicstyle=\ttfamily\small,
  backgroundcolor=\color{codebg},
  frame=single, rulecolor=\color{codeframe},
  keywordstyle=\color{thmcolor}\bfseries,
  commentstyle=\color{remcolor}\itshape,
  stringstyle=\color{defcolor},
  showstringspaces=false, breaklines=true,
  columns=fullflexible, keepspaces=true,
  xleftmargin=4pt, xrightmargin=4pt, aboveskip=8pt, belowskip=8pt
}
\lstnewenvironment{python}[1][]{\lstset{language=Python, #1}}{}
\lstnewenvironment{rcode}[1][]{\lstset{language=R, #1}}{}

% ---- Chapter header: which courses it draws on and what it builds on -------
% Usage, directly after \chapter{...}:
%   \chapterinfo{STAT2003, STAT1301}{\cref{ch:probability-spaces}}
\newcommand{\chapterinfo}[2]{%
  \begin{tcolorbox}[bookbox, before skip=2pt, colback=remcolor!6, borderline west={1pt}{0pt}{remcolor}]
    \small\textbf{Courses:} #1\par
    \textbf{Builds on:} #2
  \end{tcolorbox}}
